\subsection{等变嵌入定理}

本节中假设 \(r \neq \omega\)。

引理 3. 设 G 是一个在拓扑空间 X 上连续作用的紧拓扑群；设 A 是 X 的子集，在 G 的作用下稳定，W 是 A 的一个邻域。则存在 A 的一个开邻域 V，在 G 的作用下稳定且包含于 W 中。

置 \(\mathrm{F} = \mathrm{X} - \mathrm{W}\) 和 \(\mathrm{V} = \mathrm{X} - \mathrm{GF}\)。则 \(\mathrm{V}\) 是开的（General Topology, Chap. III, §4, no. 1, Cor. 1 of Prop. 1），在 \(\mathrm{G}\) 的作用下稳定，且 \(\mathrm{A} \subset \mathrm{V} \subset \mathrm{W}\)。

定理 1. 设 G 是紧李群，\((g,x)\mapsto gx\) 是 G 在 X 上的 \(C^{r}\) 类左作用律，A 是 X 的紧子集。存在 G 在有限维向量空间 E 上的解析线性表示 \(\rho\)、一个态射 \(\varphi:X\to E\)（为 \(C^{r}\) 类，与 G 的作用相容）以及 A 的一个开邻域 U，在 G 的作用下稳定，使得 \(\varphi\) 在 U 上的限制是嵌入。

用紧集 GA 替换 A 后，可归结为 A 在 G 的作用下稳定的情形。

设 \(E_0\) 是一个有限维向量空间，使得 \(\mathcal{C}^r(\mathrm{X};\mathrm{E}_0)\) 中存在在 A 的邻域中为嵌入的元素（no. 1, Prop. 4）；具有此性质的态射所成的集合 \(\mathcal{P}\) 是 \(\mathcal{C}^r(\mathrm{X};\mathrm{E}_0)\) 的非空开子集（no. 1, Prop. 3）。考虑紧群 G 在空间 \(\mathcal{C}^r(\mathrm{X};\mathrm{E}_0)\) 上的连续线性表示（§6, no. 4, Lemma 4）。由 Peter–Weyl 定理（Spectral Theory, in preparation），在 G 的作用下稳定的有限维子空间之并在 \(\mathcal{C}^r(\mathrm{X};\mathrm{E}_0)\) 中稠密；因此，存在元素 \(\varphi_0\) 属于 \(\mathcal{P}\)，使得映射 \(x \mapsto \varphi_0(gx)\)（\(g \in G\)）生成有限维向量子空间 \(E_1\)，它是 \(\mathcal{C}^r(\mathrm{X};\mathrm{E}_0)\) 的子空间，该子空间显然在 G 的作用下稳定。

取 E 为空间 \(\operatorname{Hom}_{\mathbf{R}}(\mathrm{E}_{1},\mathrm{E}_{0})\)，取 \(\rho\) 为 G 在 E 上由其在 \(E_{1}\) 上的作用诱导的表示，并令 \(\varphi:X\to E\) 为将 \(x\in X\) 赋予线性映射 \(\psi\mapsto\psi(x)\)（从 \(E_{1}\) 到 \(E_{0}\)）的映射。这是 \(C^{r}\) 类态射；对于 \(x\in X\)、\(g\in G\)、\(\psi\in E_{1}\)，有（记 \(\tau(g)\) 为 X 的自同构 \(x\mapsto gx\)）：

\[
\varphi (g x) (\psi) - \psi (g x) = \varphi (x) (\psi \circ \tau (g)) = (\rho (g) \varphi (x)) (\psi).
\]

令 \(\alpha : \mathrm{Hom}_{\mathbf{R}}(\mathrm{E}_1, \mathrm{E}_0) \to \mathrm{E}_0\) 为线性映射 \(u \mapsto u(\varphi_0)\)；有 \(\alpha \circ \varphi = \varphi_0\)，所以 \(\varphi\) 在 A 的邻域中是嵌入（因为 \(\varphi_0\) 是嵌入）。因此存在 A 的一个开邻域 U，使得 \(\varphi\) 在 U 上的限制是嵌入；由引理 3 可将 U 选为在 G 的作用下稳定的，定理得证。

推论 1. 假设 X 是紧的。存在 G 在有限维向量空间 E 上的解析线性表示 \(\rho\) 以及一个嵌入 \(\varphi: X \to E\)，使得 \(\varphi(gx) = \rho(g)\varphi(x)\) 对 \(g \in G, x \in X\) 成立。

推论 2. 设 H 是 G 的闭子群。存在 G 在有限维向量空间 E 上的解析线性表示以及一点 \(v \in E\)，其固定子群为 H。

将推论 1 应用于 G 在紧流形 G/H 上的典则作用。这给出解析线性表示 \(\rho: \mathrm{G} \to \mathbf{GL}(\mathrm{E})\) 以及嵌入 \(\varphi: \mathrm{G} / \mathrm{H} \to \mathrm{E}\)，使得 \(\varphi(gx) = \rho(g)\varphi(x)\)，其中 \(g \in \mathrm{G}, x \in \mathrm{G} / \mathrm{H}\)。令 \(\bar{e} \in \mathrm{G} / \mathrm{H}\) 为 \(e \in \mathrm{G}\) 的等价类，并令 \(v = \varphi(\bar{e})\) 为其像。对于所有 \(g \in \mathrm{G}\)，有 \(\rho(g)v = v \Longleftrightarrow \varphi(g\bar{e}) = \varphi(\bar{e}) \Longleftrightarrow g\bar{e} = \bar{e} \Longleftrightarrow g \in \mathrm{H}\)。

推论 3. 假设 X 是仿紧的。存在实希尔伯特空间 E、G 在 E 上的连续酉表示 \({}^{10}\) \(\rho\) 以及一个嵌入 \(\varphi: X \to E\)（为 \(C^{r}\) 类），使得 \(\varphi(gx) = \rho(g)\varphi(x)\) 对所有 \(g \in G\) 和所有 \(x \in X\) 成立。

空间 X/G 是局部紧的（General Topology, Chap. III, §4, no. 5, Prop. 11）。它的连通分支是 X 的连通分支的像，而 X 的连通分支在无穷远处可数（General Topology, Chap. I, §9, no. 10, Th. 5）；因此这些像自身也在无穷远处可数，从而 X/G 是仿紧的（loc. cit.）。所以存在 X/G 的一个局部有限覆盖 \((\mathrm{U}_{\alpha}^{\prime})_{\alpha\in\mathrm{I}}\)，其成员为相对紧开集；还存在覆盖 \((\mathrm{V}_{\alpha}^{\prime})_{\alpha\in\mathrm{I}}\)，使得 \(\overline{V}_{\alpha}^{\prime}\subset U_{\alpha}^{\prime}\) 对所有 \(\alpha\in I\) 成立（General Topology, Chap. IX, §4, no. 3, Cor. 1 of Th. 3）。取逆像后，得到 X 的两个局部有限覆盖 \((\mathrm{U}_{\alpha})_{\alpha\in\mathrm{I}}\) 和 \((\mathrm{V}_{\alpha})_{\alpha\in\mathrm{I}}\)，其成员为在 G 下稳定的相对紧开集，并且 \(\overline{V}_{\alpha}\subset U_{\alpha}\) 对所有 \(\alpha\in I\) 成立。

对于所有 \(\alpha \in \mathrm{I}\)，存在一个表示 \(\rho_{\alpha}\)，它作用在有限维实向量空间 \(\mathrm{E}_{\alpha}\) 上，以及一个与 G 的作用相容的态射 \(\varphi_{\alpha} \in \mathcal{C}^r (\mathrm{X};\mathrm{E}_{\alpha})\)，且其在 \(\mathrm{U}_{\alpha}\) 上的限制是嵌入（定理 1）。对于 \(\alpha \in \mathrm{I}\)，取一个 \(a_{\alpha}\)，它是 X 上的 \(\mathrm{C}^r\) 类数值函数，在 \(\mathrm{V}_{\alpha}\) 上等于 1，在 \(\mathrm{U}_{\alpha}\) 外等于 0（Differentiable and Analytic Manifolds, Results, 5.3.6）。置 \(b_{\alpha}(x) = \int_{\mathrm{G}} a_{\alpha}(gx) dg\) 对于 \(x \in \mathrm{X}\) 成立。函数 \(b_{\alpha}\) 是 \(\mathrm{C}^r\) 类的，在 G 下不变（§6, no. 4, Cor. 2），在 \(\mathrm{V}_{\alpha}\) 上等于 1，在 \(\mathrm{U}_{\alpha}\) 外等于 0。在每个 \(\mathrm{E}_{\alpha}\) 上取一个 G-不变的希尔伯特内积（§1, no. 1），并在 R 上取其典则希尔伯特结构；令 E 为族 \((\mathrm{E}_{\alpha} \oplus \mathbf{R})_{\alpha \in \mathrm{I}}\) 的希尔伯特和，并令 \(\rho\) 为 G 在 E 上由各 \(\rho_{\alpha}\) 以及 G 在 R 上的平凡作用诱导的表示。对于 \(x \in \mathrm{X}\)，置 \(\varphi(x) = (b_{\alpha}(x)\varphi_{\alpha}(x), b_{\alpha}(x))_{\alpha \in \mathrm{I}}\)。则 \(\varphi\) 是从 X 到 E 的 \(\mathrm{C}^r\) 类态射，与 G 的作用相容；如命题 4 的证明（no. 1）中那样验证 \(\varphi\) 是嵌入，即得该推论。

